\chapter{Conservation généalogique de l'information, entropies de réduction et structure thermique}
\label{chap:informacion-entropia-landauer}

L'information que transporte un état HMT ne se réduit pas à la coordonnée que
publie l'un de ses lecteurs. La coordonnée visible, le feuillet, l'orientation,
le quotient, la retenue, le chemin et le sceau de clôture appartiennent au même
état enrichi, bien qu'une évaluation concrète n'en montre qu'une partie.
Cette distinction permet de situer dans un cadre unique trois opérations que la
description macroscopique présente habituellement séparément : conserver une histoire,
la réduire à une lecture et attribuer une énergie à l'information effectivement
effacée.

Le chapitre développe cette chaîne sans identifier ses échelons. On
formule d'abord la mémoire de la fibre et le critère exact de récupérabilité. Ensuite,
on démontre la nullité du terme logique de Landauer pour une mise à jour
bijective de l'état enrichi du TPK. On distingue alors cinq
entropies qui possèdent des domaines et des fonctions différents. Enfin, le quotient
aréal--volumétrique et l'horloge de \(108\) pas fournissent la structure qui
relie information triadique, action et température.

\section{Désintégration de la lecture visible sur les fibres et récupération de la mémoire généalogique}

Soit \(\Omega\) un espace fini d'états enrichis et soit
\[
 q:\Omega\longrightarrow Y
\]
un lecteur visible. La valeur \(q(x)\) est une coordonnée de l'état \(x\), non sa
description exhaustive. Notons
\[
 M:\Omega\longrightarrow\mathcal M
\]
le registre conjoint des coordonnées généalogiques que le lecteur ne publie pas :
feuillet, orientation, quotient ou retenue, phase, chemin et sceau, selon le domaine
considéré.

\begin{definition}[Lecteur enrichi et mémoire de fibre]
Le lecteur enrichi associé à \(q\) et \(M\) est
\begin{equation}
 \widehat q=(q,M):\Omega\longrightarrow Y\times\mathcal M.
 \label{eq:lector-enriquecido-informacion}
\end{equation}
La mémoire de fibre de \(q\) est l'information nécessaire pour distinguer les
états contenus dans une même fibre \(q^{-1}(y)\).
\end{definition}

L'extension résidu--quotient d'APP est l'exemple arithmétique élémentaire. Dans
chaque cellule, les deux divisions euclidiennes sont conservées
\[
 i+j=9q^+_{ij}+R^+_{ij},
 \qquad
 ij=9q^\times_{ij}+R^\times_{ij}.
\]
La réduction modulo neuf publie \(R^+\) ou \(R^\times\) ; le relèvement
conserve simultanément \((R,q)\). En sommant sur le domaine nonadique, on
obtient l'identité exacte
\begin{equation}
 2025-810
 =54+9\cdot129.
 \label{eq:app-residuo-cociente-informacion}
\end{equation}
Le premier terme à droite est la différence entre les résidus visibles ;
le deuxième recompose la différence entre les quotients. L'égalité démontre
que la projection résiduelle ne détermine pas à elle seule l'opération entière.

Soit maintenant \(X\) une variable aléatoire à valeurs dans \(\Omega\). Nous travaillerons
avec des entropies adimensionnelles en nats. La règle de la chaîne donne
\begin{equation}
 \mathsf H(X)
 =\mathsf H(q(X))+\mathsf H(X\mid q(X)).
 \label{eq:cadena-entropia-fibra}
\end{equation}
Le deuxième terme quantifie l'information généalogique qui subsiste à l'intérieur
de la fibre de la lecture.

\begin{theorem}[Récupérabilité par mémoire de fibre]
\label{thm:recuperabilidad-memoria-fibra}
Si le lecteur enrichi \(\widehat q=(q,M)\) est injectif sur le support de
\(X\), alors
\begin{equation}
 \boxed{\mathsf H(X\mid q(X),M(X))=0.}
 \label{eq:entropia-condicional-memoria-cero}
\end{equation}
L'entropie \(\mathsf H(X\mid q(X))\) mesure exactement l'information qui serait
perdue en conservant la coordonnée visible et en écartant la mémoire.
\end{theorem}

\begin{proof}
L'injectivité de \(\widehat q\) fournit une inverse sur son image. Par
conséquent, la paire \((q(X),M(X))\) détermine \(X\) et l'entropie conditionnelle est
nulle. Si l'on omet \(M(X)\), les entrées d'une même fibre sont
identifiées ; la règle de la chaîne \eqref{eq:cadena-entropia-fibra} attribue à
cette identification le terme \(\mathsf H(X\mid q(X))\).
\end{proof}

La définition généalogique du nombre s'insère précisément ici. Un nombre
HMT est une section compatible de cylindres survivants, accompagnée de la mémoire
de son chemin, de son orientation, de sa signature et de sa clôture. Sa valeur réelle est l'évaluation de cette
section. L'évaluation peut oublier des coordonnées ; la section les conserve.

\section{Réversibilité logique et terme de Landauer}

Soit
\[
 U:\Omega\longrightarrow\Omega
\]
une mise à jour de l'état enrichi du TPK. La question de Landauer se
formule sur l'application complète, avant d'appliquer un lecteur visible ou de réinitialiser
un registre ; sa formulation thermodynamique de l'effacement logique remonte à
Landauer \cite{Landauer1961}.

\begin{theorem}[Landauer nul pour le transport complet]
\label{thm:landauer-nulo-tpk-anexo}
Si \(U\) est bijective, alors, pour toute distribution d'entrée \(X\),
\begin{equation}
 \mathsf H(X\mid U(X))=0.
 \label{eq:landauer-entropia-cero-anexo}
\end{equation}
Dans le régime idéal de Landauer, le terme minimal associé exclusivement à
l'effacement logique interne est
\begin{equation}
 \boxed{
 Q_{\rm L}^{\min}[U]
 =k_B\Theta\,\mathsf H(X\mid U(X))=0.}
 \label{eq:landauer-nulo-anexo}
\end{equation}
\end{theorem}

\begin{proof}
La bijection fournit \(U^{-1}\), de sorte que
\(X=U^{-1}(U(X))\). La sortie complète détermine l'entrée et l'entropie
conditionnelle s'annule. La forme informationnelle du principe de Landauer attribue
le coût d'effacement au terme \(k_B\Theta\mathsf H(X\mid U(X))\), qui, par
l'égalité précédente, vaut zéro.
\end{proof}

Le résultat localise la conservation dans l'espace enrichi. Pour une
sortie projetée \(q(U(X))\), l'information non publiée est
\begin{equation}
 \mathsf H\bigl(X\mid q(U(X))\bigr).
 \label{eq:informacion-proyeccion-anexo}
\end{equation}
Une réinitialisation physique qui identifie des états distincts doit transférer cette
information à l'environnement. En particulier, la réinitialisation d'un TRIT uniforme vers un unique
état possède
\[
 \mathsf H(X)=\log3,
 \qquad
 \boxed{Q_{\rm reset}\geq k_B\Theta\log3.}
\]
La mise à jour réversible, la projection et la réinitialisation sont ainsi des opérations
distinctes. La première conserve la généalogie ; la deuxième réduit la lecture ; la
troisième élimine physiquement un registre. La dissipation supplémentaire due au
contrôle, au bruit ou à une vitesse finie appartient à la réalisation du
dispositif et non au terme logique d'effacement de
\eqref{eq:landauer-nulo-anexo}.

\section{\(5\) notions d'entropie et leurs domaines}

Le mot \emph{entropie} désigne dans cette construction cinq fonctionnelles
liées, mais non interchangeables. Leur séparation évite d'attribuer une température
à un dénombrement purement combinatoire ou une perte d'information à une évolution
qui ne fait que changer de lecture.

\subsection{Entropie conditionnelle des fibres du lecteur visible et perte de généalogie}

La quantité
\[
 \mathsf H_{\rm fib}(q;X):=\mathsf H(X\mid q(X))
\]
mesure l'information de la généalogie qu'une lecture visible ne distingue pas.
Son domaine est un lecteur et une distribution sur l'état enrichi. Elle
s'annule lorsque le lecteur est fidèle sur le support considéré.

\subsection{Entropie informationnelle d'une décision triadique}

Pour une distribution \(p=(p_{-1},p_0,p_{+1})\) sur les trois états
visibles du TRIT,
\[
 I_{\rm TRIT}(p)=-\sum_{a\in\{-1,0,+1\}}p_a\log p_a.
\]
La distribution uniforme produit
\begin{equation}
 \boxed{I_{\rm TRIT}=\log3.}
 \label{eq:informacion-trit-log3-anexo}
\end{equation}
Cette quantité mesure une décision locale entre trois alternatives visibles.

\subsection{Entropie topologique des histoires admissibles}

Si \(A\) est la matrice d'adjacence d'un sous-shift fini, son entropie
topologique est
\[
 h_{\rm top}(A)=\log\rho(A),
\]
où \(\rho(A)\) est le rayon spectral. Cette fonctionnelle mesure le taux
asymptotique de croissance des histoires compatibles, non l'incertitude
d'une décision isolée.

\subsection{Entropie combinatoire et entropie statistique maximale}

Pour des niveaux d'énergie \(\varepsilon_i\) de dégénérescences \(g_i\), une
occupation fermionique \(n_i\in\{0,\ldots,g_i\}\) possède
\[
 \Omega(g_i,n_i)=\binom{g_i}{n_i}
\]
microétats. Avec un nombre et une énergie totaux fixés,
\begin{equation}
 \Omega(N,E)=
 \sum_{\substack{(n_i):\ \sum_i n_i=N\\
                         \sum_i\varepsilon_i n_i=E}}
 \prod_i\binom{g_i}{n_i},
 \qquad
 S_{\rm mc}=\log\Omega(N,E).
 \label{eq:entropia-microcanonica-anexo}
\end{equation}
En fixant le nombre et l'énergie comme valeurs moyennes, la maximisation de Shannon
\cite{Shannon1948} donne
la distribution produit
\[
 \mathbb P(n_1,\ldots,n_M)
 \propto
 \exp\!\left[-\beta\sum_j\varepsilon_jn_j
              +\beta\mu\sum_j n_j\right]
\]
et, par conséquent,
\begin{equation}
 \boxed{
 \langle n_j\rangle_{\rm FD}
 =\frac{1}{e^{\beta(\varepsilon_j-\mu)}+1}.}
 \label{eq:fermi-dirac-entropia-anexo}
\end{equation}
L'occupation bosonique remplace le dénombrement binomial par
\(\binom{g+n-1}{n}\) et produit
\[
 \langle n_j\rangle_{\rm BE}
 =\frac{1}{e^{\beta(\varepsilon_j-\mu)}-1}.
\]
Dans les deux cas, \(\beta\) et \(\mu\) sont déterminés par les contraintes
macroscopiques déclarées.

\subsection{Entropie de von Neumann de \(1\) état réduit}

Pour une matrice de densité \(\rho_A\) sur une région ou une sous-algèbre,
\[
 S_{\rm vN}(\rho_A)=-\operatorname{Tr}(\rho_A\log\rho_A).
\]
Son domaine est un état réduit. Lorsque \(\rho_A\) provient d'une
purification globale, cette quantité mesure l'intrication à travers la coupure.
Le chapitre suivant construit la réduction triadique, le Hamiltonien modulaire
et la quantification d'aire qui lui donnent une résidence géométrique précise.

\begin{proposition}[Séparation typée des cinq entropies]
\label{prop:cinco-entropias-tipadas}
Les quantités \(\mathsf H_{\rm fib}\), \(I_{\rm TRIT}\), \(h_{\rm top}\),
\(S_{\rm mc}\) et \(S_{\rm vN}\) ne se substituent pas les unes aux autres : elles dépendent,
respectivement, d'une fibre de lecture, d'une distribution locale triadique, d'un
système d'histoires, d'un ensemble de microétats compatible avec des contraintes
et d'un état réduit. Une application physique ne peut les relier qu'après avoir
déclaré le coarse--graining, l'ensemble, le Hamiltonien ou la purification
qui fait passer d'un domaine au suivant.
\end{proposition}

\begin{proof}
Chaque fonctionnelle est définie sur un type d'entrée différent. L'entropie
conditionnelle nécessite un lecteur ; \(I_{\rm TRIT}\), une distribution sur trois
symboles ; \(h_{\rm top}\), une dynamique de mots ; \(S_{\rm mc}\), une
contrainte macrocanonique ; et \(S_{\rm vN}\), un opérateur de densité. Une
égalité numérique accidentelle ne fournit pas les applications qui changent ces types.
\end{proof}

\section{Incidence surfacique--volumétrique et entropie topologique}

Le calendrier du TPK induit la charnière entre les modes additif et
multiplicatif. Dans le bloc primitif
\[
 \tau_3=(+1,-1,0),
\]
la règle
\[
 +1:m\mapsto\Sigma,
 \qquad
 -1:m\mapsto\Pi,
 \qquad
 0:m\mapsto m
\]
produit l’orbite cyclique \((\Sigma,\Pi,\Pi)\). Le lecteur géométrique
\[
 \Sigma\longmapsto\mathscr H_{\rm ar},
 \qquad
 \Pi\longmapsto\mathscr H_{\rm vol}
\]
transporte ces modes vers les feuillets surfacique et volumétrique.

\begin{theorem}[Incidence nécessaire des feuillets]
\label{thm:incidencia-hojas-anexo}
Dans l’ordre \((\mathscr H_{\rm ar},\mathscr H_{\rm vol})\), la matrice
primitive des incidences est
\begin{equation}
 \boxed{
 F_{\rm av}=
 \begin{pmatrix}0&1\\1&1\end{pmatrix}.}
 \label{eq:fav-anexo-informacion}
\end{equation}
Les calendriers de longueurs \(9,27,108\) produisent
\[
 N_9=3F_{\rm av},
 \qquad N_{27}=9F_{\rm av},
 \qquad N_{108}=36F_{\rm av}.
\]
\end{theorem}

\begin{proof}
Les paires consécutives du cycle sont
\(\Sigma\to\Pi\), \(\Pi\to\Pi\) et \(\Pi\to\Sigma\). Chacune apparaît une
fois et \(\Sigma\to\Sigma\) n’apparaît pas. Cela fixe les quatre entrées de
\(F_{\rm av}\). Les calendriers longs contiennent, respectivement, trois,
neuf et trente-six copies du bloc primitif.
\end{proof}

\begin{theorem}[Loi de Fibonacci et mesure d’entropie maximale]
\label{thm:fibonacci-parry-anexo}
Pour \(n\geq1\),
\begin{equation}
 F_{\rm av}^{n}
 =\begin{pmatrix}F_{n-1}&F_n\\F_n&F_{n+1}\end{pmatrix}.
 \label{eq:potencias-fav-anexo}
\end{equation}
Le sous-décalage à mémoire un défini par \(F_{\rm av}\) a
\begin{equation}
 \boxed{h_{\rm top}=\log\varphi},
 \qquad
 \boxed{\frac{\mu_{\rm ar}}{\mu_{\rm vol}}=\varphi^{-2}}.
 \label{eq:entropia-parry-hojas-anexo}
\end{equation}
\end{theorem}

\begin{proof}
L’identité matricielle s’obtient par récurrence à partir de
\(F_{n+1}=F_n+F_{n-1}\). La racine de Perron de \(F_{\rm av}\) est \(\varphi\),
d’où \(h_{\rm top}=\log\varphi\). Comme la matrice est symétrique, les poids
de Parry \cite{Parry1964} sont proportionnels au carré d’un vecteur de Perron
\((1,\varphi)\), c’est-à-dire à \((1,\varphi^2)\).
\end{proof}

Les expressions \(\log3\) et \(\log\varphi\) sont ainsi séparées par leur
généalogie. La première quantifie un choix local uniforme entre trois
états ; la seconde mesure la croissance des histoires permises par
l’incidence des feuillets. La relation surfacique--volumétrique est une propriété du
retour complet et apparaît donc avec l’exposant deux.

\section{Transduction de Boltzmann et énergie par décision triadique}

Soient \(\mathcal L_T\) la droite de durée, \(\mathcal L_\Theta\) la droite de
température, \(\mathcal L_S\) la droite d'action et \(\mathcal L_E\) la droite
d'énergie. La constante de Boltzmann est l'isomorphisme positif
\begin{equation}
 \boxed{
 k_B:\mathcal L_\Theta\longrightarrow\mathcal L_E,
 \qquad [k_B]=\mathbf E-\boldsymbol\Theta.}
 \label{eq:boltzmann-transductor-anexo}
\end{equation}
Ainsi, \(\beta=(k_B\Theta)^{-1}\) appartient à la droite inverse de l'énergie.

Le lecteur déterminantal de l'action construit
\[
 \hbar_{\rm int}
 =\eta_{\rm ret}\,10^{-34}\mathcal U_S\in\mathcal L_S.
\]
Si \(t_0\in\mathcal L_T^+\) est la durée interne élémentaire, l'horloge de
retour complet fixe
\[
 \omega_0=\frac{2\pi}{108t_0},
 \qquad
 \varepsilon_{\rm clk}=\hbar_{\rm int}\omega_0\in\mathcal L_E.
\]

\begin{definition}[Échelle thermique triadique]
La section \(\Theta_{\rm clk}\in\mathcal L_\Theta^+\) est définie par
\begin{equation}
 \boxed{
 k_B\Theta_{\rm clk}\log3
 =\hbar_{\rm int}\frac{2\pi}{108t_0}.}
 \label{eq:boltzmann-accion-informacion-anexo}
\end{equation}
\end{definition}

L'égalité fixe l'énergie par décision triadique et relie trois objets déjà
construits : l'information locale \(\log3\), l'action interne et la fréquence
de l'horloge. Le produit \(k_B\Theta_{\rm clk}\) est déterminé dans la
carte interne. Le choix des bases thermique et énergétique sépare ensuite la
coordonnée de \(k_B\) et celle de \(\Theta_{\rm clk}\). Dans la carte kelvin--joule,
la définition métrologique publie
\[
 k_B^{\rm SI}=1.380649\times10^{-23}\ \mathrm{J\,K^{-1}}.
\]

\section{Opérateur thermique sur les histoires}

Soient \(\mathcal E\) l'ensemble des arêtes d'un graphe fini d'extensions,
\(a_*(e)\geq0\) l'incrément adimensionnel d'action et
\(\epsilon_*\in\mathcal L_E^+\) une section positive d'énergie. Pour
\(\beta\in\mathcal L_E^{-1}\), nous définissons
\begin{equation}
 (\mathcal L_{\beta,*}f)(y)
 =\sum_{e:x\to y}
 \exp[-\beta\epsilon_*a_*(e)]f(x).
 \label{eq:operador-transferencia-termica-anexo}
\end{equation}
L'exposant est adimensionnel et la positivité de l'opérateur conserve la
composition des histoires. La normalisation de pression sélectionne l'état
thermique au moyen de
\begin{equation}
 \boxed{\rho(\mathcal L_{\beta,*})=1.}
 \label{eq:presion-espectral-termica-anexo}
\end{equation}

\begin{proposition}[Confluence de la branche informationnelle et de la branche spectrale]
\label{prop:confluencia-termica-anexo}
Une réalisation thermique HMT--MD est déterminée lorsqu'une même valeur de
\(\beta=(k_B\Theta)^{-1}\) satisfait la normalisation spectrale
\eqref{eq:presion-espectral-termica-anexo} et l'échelle action--information
\eqref{eq:boltzmann-accion-informacion-anexo}. La première équation pondère les
histoires ; la deuxième fixe l'énergie de la décision locale. Leur confluence
transforme la comptabilité généalogique en un ensemble thermique sans identifier
information, entropie et énergie.
\end{proposition}

La thermodynamique apparaît donc après la conservation généalogique.
L'état complet conserve l'histoire ; le lecteur définit le coarse--graining ;
le dénombrement organise les microétats ; et le transducteur de Boltzmann attribue une
échelle énergétique à l'information sous une température déclarée.
